\chapter{Implementation and Mathematical Foundations}

\section{Working Example}

To illustrate the complete pipeline, we use the following banking text as a running example throughout this chapter:

\begin{quote}
\textit{``Account Holder: Sai Manoj Yadav, Account Number: 123456789012, IFSC Code: SBIN0001234, Phone: 9876543210, Email: saimanoj@example.com, Transfer Amount: Rs.~45,000''}
\end{quote}

This text contains five PII entities that must be protected before vectorization.

\section{Threat Detection Engine}

The threat detection engine (\texttt{security.py}) scans all input text against two categories of attack patterns using precompiled regular expressions.

\subsection{Attack Pattern Categories}

\begin{table}[H]
\centering
\caption{Threat detection regex patterns}
\label{tab:threats}
\begin{tabular}{lll}
\toprule
\textbf{Category} & \textbf{Pattern Example} & \textbf{Attack Type} \\
\midrule
Code Injection & \texttt{<script>.*</script>} & XSS \\
Code Injection & \texttt{DROP\textbackslash s+TABLE} & SQL Injection \\
Code Injection & \texttt{rm\textbackslash s+-rf} & Shell Injection \\
Prompt Injection & \texttt{ignore.*previous.*instructions} & LLM Jailbreak \\
Prompt Injection & \texttt{you are now.*DAN} & Role Manipulation \\
Prompt Injection & \texttt{disregard.*instructions} & Context Hijacking \\
\bottomrule
\end{tabular}
\end{table}

\subsection{Detection Algorithm}

\begin{algorithm}[H]
\caption{Threat Detection}
\begin{algorithmic}[1]
\Procedure{DetectThreat}{$\text{text}$}
    \State $\text{text\_lower} \gets \text{text.lower()}$
    \For{each $\text{pattern}$ in $\text{CODE\_INJECTION\_PATTERNS}$}
        \If{$\text{regex.search(pattern, text\_lower)}$}
            \State \Return $(\texttt{True},\ \texttt{``CODE\_INJECTION''})$
        \EndIf
    \EndFor
    \For{each $\text{pattern}$ in $\text{PROMPT\_INJECTION\_PATTERNS}$}
        \If{$\text{regex.search(pattern, text\_lower)}$}
            \State \Return $(\texttt{True},\ \texttt{``PROMPT\_INJECTION''})$
        \EndIf
    \EndFor
    \State \Return $(\texttt{False},\ \texttt{None})$
\EndProcedure
\end{algorithmic}
\end{algorithm}

\section{PII Detection: Context Lookbehind Algorithm}

\subsection{The Disambiguation Problem}

Indian banking documents frequently contain digit strings of identical length that represent different entity types. For example, both \texttt{9876543210} (phone number) and \texttt{9876543210} (account number) are 10-digit sequences. The system must determine the correct classification.

\subsection{25-Character Context Lookbehind}

For every number matched by the regex \texttt{\textbackslash b\textbackslash d\{9,18\}\textbackslash b}, the algorithm examines the 25 characters immediately preceding the match:

\begin{equation}
\text{context} = \text{text}[\max(0,\ p_{\text{start}} - 25)\ :\ p_{\text{start}}].\text{lower()}
\end{equation}

where $p_{\text{start}}$ is the starting position of the matched digit sequence.

\begin{algorithm}[H]
\caption{PII Classification with Context Lookbehind}
\begin{algorithmic}[1]
\Procedure{ClassifyNumber}{$\text{text},\ p_{\text{start}},\ p_{\text{end}},\ \text{number}$}
    \State $\text{context} \gets \text{text}[\max(0, p_{\text{start}}-25) : p_{\text{start}}].\text{lower()}$
    \State $\text{length} \gets p_{\text{end}} - p_{\text{start}}$
    \If{``phone'' $\in$ context \textbf{or} ``mob'' $\in$ context \textbf{or} ``cell'' $\in$ context}
        \State \Return \texttt{PHONE}
    \ElsIf{``acc'' $\in$ context \textbf{or} ``account'' $\in$ context}
        \State \Return \texttt{ACCOUNT}
    \ElsIf{$\text{length} = 10$ \textbf{and} number[0] $\in$ \{6,7,8,9\}}
        \State \Return \texttt{PHONE} \Comment{Indian mobile numbers start with 6--9}
    \Else
        \State \Return \texttt{ACCOUNT}
    \EndIf
\EndProcedure
\end{algorithmic}
\end{algorithm}

\subsection{Applied to Our Example}

\begin{table}[H]
\centering
\caption{PII detection results for running example}
\label{tab:pii_results}
\begin{tabular}{llll}
\toprule
\textbf{Value} & \textbf{Context (25 chars before)} & \textbf{Classification} \\
\midrule
\texttt{123456789012} & \texttt{``yadav, account number: ''} & ACCOUNT \\
\texttt{SBIN0001234} & Regex: \texttt{[A-Z]\{4\}0[A-Z0-9]\{6\}} & IFSC \\
\texttt{9876543210} & \texttt{``0001234, phone: ''} & PHONE \\
\texttt{saimanoj@example.com} & Regex: \texttt{\textbackslash S+@\textbackslash S+} & EMAIL \\
\bottomrule
\end{tabular}
\end{table}

\section{AES-256-CBC Pseudonymization}

\subsection{Token Generation}

Each detected PII entity is replaced with a UUID-based token:

\begin{equation}
\text{token} = \text{CATEGORY\_PREFIX} + \text{``\_''} + \text{UUID4}()[:6]
\end{equation}

For example: $\texttt{123456789012} \rightarrow \texttt{ACCOUNT\_f7ce}$

\subsection{AES-256-CBC Encryption of Mappings}

The mapping between tokens and original values is encrypted before storage in MongoDB using AES-256 in CBC mode.

\begin{definition}[AES-256-CBC Encryption]
Given plaintext $P$, a 256-bit key $K$, and a random 128-bit initialization vector $IV$:
\begin{align}
C_0 &= IV \\
C_i &= E_K(P_i \oplus C_{i-1}), \quad i = 1, 2, \ldots, n
\end{align}
where $E_K$ is the AES block cipher with key $K$, $P_i$ is the $i$-th plaintext block (16 bytes with PKCS7 padding), and $\oplus$ denotes XOR.
\end{definition}

\textbf{Key derivation:}
\begin{equation}
K = \text{SHA-256}(\texttt{SECRET\_KEY})
\end{equation}

The key is derived from the environment variable \texttt{SECRET\_KEY} using SHA-256, producing exactly 32 bytes (256 bits). The key never appears in the database.

\textbf{Storage format} in MongoDB:
\begin{equation}
\text{stored\_value} = \text{Base64}(IV \| \text{Ciphertext})
\end{equation}

\subsection{Decryption (Controlled De-pseudonymization)}

De-pseudonymization occurs only when:
\begin{enumerate}
    \item The requester is authenticated via JWT.
    \item The \texttt{owner\_id} in the JWT matches the \texttt{owner\_id} in the pseudonym mapping.
\end{enumerate}

\begin{equation}
P_i = D_K(C_i) \oplus C_{i-1}, \quad \text{where } C_0 = IV
\end{equation}

\section{Text-to-Vector Conversion}

\subsection{Tokenization}

The pseudonymized text is tokenized into sub-word units using the model's vocabulary $V$ of size $|V| \approx 30,000$:

\begin{equation}
\text{token\_ids} = \text{Tokenizer}(\text{``Account Number: ACCOUNT\_f7ce ...''})
\end{equation}

Each word $w$ is mapped to an integer ID: $w \mapsto \text{id} \in \{0, 1, \ldots, |V|-1\}$.

\subsection{Embedding Lookup}

Each token ID is mapped to a $d$-dimensional vector (where $d = 3072$ for Gemini Embedding) through an embedding matrix $\mathbf{E} \in \mathbb{R}^{|V| \times d}$:

\begin{equation}
\mathbf{e}_i = \mathbf{E}[\text{token\_id}_i] \in \mathbb{R}^{3072}
\end{equation}

\subsection{Self-Attention Mechanism}

The Gemini transformer processes the token embeddings using the \textbf{Scaled Dot-Product Self-Attention} mechanism:

\begin{equation}
\boxed{\text{Attention}(\mathbf{Q}, \mathbf{K}, \mathbf{V}) = \text{softmax}\left(\frac{\mathbf{Q}\mathbf{K}^T}{\sqrt{d_k}}\right) \mathbf{V}}
\label{eq:attention}
\end{equation}

where:
\begin{itemize}
    \item $\mathbf{Q} = \mathbf{X}\mathbf{W}_Q \in \mathbb{R}^{n \times d_k}$ (Query matrix)
    \item $\mathbf{K} = \mathbf{X}\mathbf{W}_K \in \mathbb{R}^{n \times d_k}$ (Key matrix)
    \item $\mathbf{V} = \mathbf{X}\mathbf{W}_V \in \mathbb{R}^{n \times d_v}$ (Value matrix)
    \item $\mathbf{X} \in \mathbb{R}^{n \times d}$ is the matrix of input embeddings
    \item $n$ is the number of tokens, $d_k = 256$ is the key dimension
\end{itemize}

The scaling factor $\frac{1}{\sqrt{d_k}}$ prevents the dot products from growing too large, which would push the softmax into regions of extremely small gradients.

\subsection{Multi-Head Attention}

The model runs $h = 12$ parallel attention heads, each learning different relational patterns:

\begin{equation}
\text{MultiHead}(\mathbf{Q}, \mathbf{K}, \mathbf{V}) = \text{Concat}(\text{head}_1, \text{head}_2, \ldots, \text{head}_{12}) \cdot \mathbf{W}_O
\end{equation}

where each head computes:
\begin{equation}
\text{head}_i = \text{Attention}(\mathbf{X}\mathbf{W}_Q^{(i)}, \mathbf{X}\mathbf{W}_K^{(i)}, \mathbf{X}\mathbf{W}_V^{(i)})
\end{equation}

\subsection{Output: Document Vector}

After all transformer layers, the model produces a single 3072-dimensional vector $\mathbf{v} \in \mathbb{R}^{3072}$ representing the entire document's semantic meaning:

\begin{equation}
\mathbf{v} = \text{GeminiEmbed}(\text{``Account Number: ACCOUNT\_f7ce ...''})\in \mathbb{R}^{3072}
\end{equation}

\section{Differential Privacy: Laplacian Noise Injection}

\subsection{Motivation: Embedding Inversion}

Without noise, an attacker with access to the stored vector $\mathbf{v}$ could train an inversion model $g_\theta$ to approximate the inverse mapping:

\begin{equation}
g_\theta(\mathbf{v}) \approx \text{original text}
\end{equation}

Song and Raghunathan~\cite{song2020information} demonstrated up to 70\% token-level reconstruction accuracy. Adding calibrated noise destroys this inverse mapping.

\subsection{The Laplace Distribution}

Noise is sampled from the Laplace distribution:

\begin{equation}
\boxed{f(x \mid \mu, b) = \frac{1}{2b} \exp\left(-\frac{|x - \mu|}{b}\right)}
\label{eq:laplace}
\end{equation}

where $\mu = 0$ (zero-centered) and $b = \Delta f / \epsilon$ (scale parameter).

\subsection{The $\epsilon$-Differential Privacy Mechanism}

For each dimension $i$ of the embedding vector:

\begin{equation}
\boxed{v'_i = v_i + \eta_i, \quad \text{where } \eta_i \sim \text{Laplace}\left(0, \frac{\Delta f}{\epsilon}\right)}
\label{eq:dp}
\end{equation}

\begin{table}[H]
\centering
\caption{Differential privacy parameters}
\label{tab:dp_params}
\begin{tabular}{ccl}
\toprule
\textbf{Parameter} & \textbf{Symbol} & \textbf{Description} \\
\midrule
Privacy budget & $\epsilon$ & Lower $\Rightarrow$ more privacy (tunable) \\
Sensitivity & $\Delta f$ & Maximum embedding shift per data point \\
Scale & $b = \Delta f / \epsilon$ & Controls noise magnitude \\
Center & $\mu$ & 0 (noise is symmetric) \\
Dimensions & $d$ & 3072 (Gemini embedding size) \\
\bottomrule
\end{tabular}
\end{table}

\subsection{Applied to Our Example}

\begin{table}[H]
\centering
\caption{Differential privacy applied to embedding dimensions}
\label{tab:dp_example}
\begin{tabular}{cccc}
\toprule
\textbf{Dimension} & \textbf{Original $v_i$} & \textbf{Noise $\eta_i$} & \textbf{Noised $v'_i$} \\
\midrule
1 & $-0.008764$ & $+0.001153$ & $-0.007611$ \\
2 & $+0.011886$ & $-0.002114$ & $+0.009772$ \\
3 & $+0.022306$ & $+0.000294$ & $+0.022600$ \\
4 & $-0.082017$ & $+0.003561$ & $-0.078456$ \\
$\vdots$ & $\vdots$ & $\vdots$ & $\vdots$ \\
3072 & $-0.019234$ & $+0.000847$ & $-0.018387$ \\
\bottomrule
\end{tabular}
\end{table}

The noise shifts each coordinate by approximately $\pm 0.003$. This is small enough to preserve cosine similarity (search accuracy) but large enough to prevent precise floating-point inversion.

\section{Vector Retrieval: Cosine Similarity}

When a user submits a query $q$, it is converted to a query vector $\mathbf{q} \in \mathbb{R}^{3072}$. The system retrieves the top-$k$ most similar document vectors using \textbf{cosine similarity}:

\begin{equation}
\boxed{\cos(\theta) = \frac{\mathbf{q} \cdot \mathbf{d}}{\|\mathbf{q}\| \times \|\mathbf{d}\|} = \frac{\sum_{i=1}^{3072} q_i \cdot d_i}{\sqrt{\sum_{i=1}^{3072} q_i^2} \times \sqrt{\sum_{i=1}^{3072} d_i^2}}}
\label{eq:cosine}
\end{equation}

\subsection{Numerical Example}

For query \textit{``What is the transfer amount?''} and our stored document vector:

\begin{align}
\mathbf{q} \cdot \mathbf{d} &= q_1 d_1 + q_2 d_2 + \cdots + q_{3072} d_{3072} = 0.8234 \\
\|\mathbf{q}\| &= \sqrt{q_1^2 + q_2^2 + \cdots + q_{3072}^2} = 0.9156 \\
\|\mathbf{d}\| &= \sqrt{d_1^2 + d_2^2 + \cdots + d_{3072}^2} = 0.9023 \\
\cos(\theta) &= \frac{0.8234}{0.9156 \times 0.9023} = \frac{0.8234}{0.8262} = 0.9966
\end{align}

Since $0.9966 \geq 0.5$ (the relevance threshold), this document is returned as a relevant context chunk.

\subsection{Relevance Threshold}

Only chunks with cosine similarity above the threshold $\tau = 0.5$ are returned:

\begin{equation}
\text{relevant\_chunks} = \{d_j \mid \cos(\mathbf{q}, \mathbf{d}_j) \geq \tau,\ \text{owner\_id}(d_j) \in \text{allowed}\}
\end{equation}

\section{Backend File Structure}

Table~\ref{tab:files} maps each backend file to its role in the security pipeline.

\begin{table}[H]
\centering
\caption{Backend file reference}
\label{tab:files}
\begin{tabular}{lll}
\toprule
\textbf{File} & \textbf{Purpose} & \textbf{Security Role} \\
\midrule
\texttt{main.py} & FastAPI orchestrator & Rate limiting, de-pseudonymization \\
\texttt{auth.py} & Registration, JWT, bcrypt & Authentication \\
\texttt{security.py} & Regex threat detection & Blocks SQLi, XSS, Prompt Injection \\
\texttt{pii\_detector.py} & PII entity extraction & Account, IFSC, Phone, Email \\
\texttt{pseudo.py} & AES-256 pseudonymization & PII protection in MongoDB \\
\texttt{rag.py} & Embedding + LLM generation & Gemini API interface \\
\texttt{dp\_embedding.py} & Laplacian noise injection & Embedding inversion defense \\
\texttt{db.py} & Pinecone CRUD & Multi-tenant vector isolation \\
\texttt{config.py} & Environment variable loader & AES key, API keys \\
\bottomrule
\end{tabular}
\end{table}
